% Options for packages loaded elsewhere
\PassOptionsToPackage{unicode}{hyperref}
\PassOptionsToPackage{hyphens}{url}
%
\documentclass[
]{article}
\usepackage{amsmath,amssymb}
\usepackage{iftex}
\ifPDFTeX
  \usepackage[T1]{fontenc}
  \usepackage[utf8]{inputenc}
  \usepackage{textcomp} % provide euro and other symbols
\else % if luatex or xetex
  \usepackage{unicode-math} % this also loads fontspec
  \defaultfontfeatures{Scale=MatchLowercase}
  \defaultfontfeatures[\rmfamily]{Ligatures=TeX,Scale=1}
\fi
\usepackage{lmodern}
\ifPDFTeX\else
  % xetex/luatex font selection
  \setmainfont[]{Inter}
\fi
% Use upquote if available, for straight quotes in verbatim environments
\IfFileExists{upquote.sty}{\usepackage{upquote}}{}
\IfFileExists{microtype.sty}{% use microtype if available
  \usepackage[]{microtype}
  \UseMicrotypeSet[protrusion]{basicmath} % disable protrusion for tt fonts
}{}
\makeatletter
\@ifundefined{KOMAClassName}{% if non-KOMA class
  \IfFileExists{parskip.sty}{%
    \usepackage{parskip}
  }{% else
    \setlength{\parindent}{0pt}
    \setlength{\parskip}{6pt plus 2pt minus 1pt}}
}{% if KOMA class
  \KOMAoptions{parskip=half}}
\makeatother
\usepackage{xcolor}
\usepackage[margin=18mm]{geometry}
\usepackage{longtable,booktabs,array}
\usepackage{calc} % for calculating minipage widths
% Correct order of tables after \paragraph or \subparagraph
\usepackage{etoolbox}
\makeatletter
\patchcmd\longtable{\par}{\if@noskipsec\mbox{}\fi\par}{}{}
\makeatother
% Allow footnotes in longtable head/foot
\IfFileExists{footnotehyper.sty}{\usepackage{footnotehyper}}{\usepackage{footnote}}
\makesavenoteenv{longtable}
\setlength{\emergencystretch}{3em} % prevent overfull lines
\providecommand{\tightlist}{%
  \setlength{\itemsep}{0pt}\setlength{\parskip}{0pt}}
\setcounter{secnumdepth}{-\maxdimen} % remove section numbering
\ifLuaTeX
  \usepackage{selnolig}  % disable illegal ligatures
\fi
\IfFileExists{bookmark.sty}{\usepackage{bookmark}}{\usepackage{hyperref}}
\IfFileExists{xurl.sty}{\usepackage{xurl}}{} % add URL line breaks if available
\urlstyle{same}
\hypersetup{
  hidelinks,
  pdfcreator={LaTeX via pandoc}}

\author{}
\date{}

\begin{document}

\hypertarget{cotizaciuxf3n-comercial--plataforma-cicuc}{%
\section{Cotización comercial --- Plataforma
CICUC}\label{cotizaciuxf3n-comercial--plataforma-cicuc}}

\textbf{Documento:} ADYAC-COT-CICUC-2026-001\\
\textbf{Fecha de emisión:} 31 de julio de 2026\\
\textbf{Vigencia de la oferta:} 15 días corridos\\
\textbf{Estado:} Borrador pendiente de completar datos tributarios y
aceptación

\hypertarget{1-proveedor}{%
\subsection{1. Proveedor}\label{1-proveedor}}

\begin{itemize}
\tightlist
\item
  \textbf{Nombre comercial:} Adyac
\item
  \textbf{Razón social:} Investigación, Tecnología y Gestión ADYAC SpA
\item
  \textbf{RUT:} 78.329.959-3
\item
  \textbf{Emisor:} ADYAC SpA
\item
  \textbf{Correo:}
  \href{mailto:hola@adyac.cl}{\nolinkurl{hola@adyac.cl}}
\item
  \textbf{Sitio web:}
  \href{https://www.adyac.cl/}{\url{https://www.adyac.cl/}}
\item
  \textbf{Domicilio:} \texttt{{[}COMPLETAR\ DOMICILIO\ TRIBUTARIO{]}}
\end{itemize}

\hypertarget{2-cliente}{%
\subsection{2. Cliente}\label{2-cliente}}

\begin{itemize}
\tightlist
\item
  \textbf{Institución solicitante:} Centro de Investigaciones Clínicas
  UC --- CICUC, Facultad de Medicina
\item
  \textbf{Razón social:} Pontificia Universidad Católica de Chile
\item
  \textbf{RUT:} 81.698.900-0
\item
  \textbf{Dirección institucional:} Av. Libertador Bernardo
  O\textquotesingle Higgins 340, Santiago
\item
  \textbf{Dirección CICUC:} Portugal 61, Santiago
\item
  \textbf{Contacto responsable:} \texttt{{[}COMPLETAR{]}}
\item
  \textbf{Correo:} \texttt{{[}COMPLETAR{]}}
\end{itemize}

\hypertarget{3-objeto-de-la-propuesta}{%
\subsection{3. Objeto de la propuesta}\label{3-objeto-de-la-propuesta}}

Adyac propone diseñar, desarrollar, validar y poner en marcha una
plataforma web modular para la gestión centralizada de estudios clínicos
oncológicos de CICUC. La solución contará con APIs documentadas para
facilitar la interacción de los equipos clínicos y una interoperabilidad
gradual con sistemas institucionales expresamente autorizados. Deberá
quedar operativa y no se limitará a un prototipo visual.

El proyecto considera descubrimiento y levantamiento, definición de una
línea base aprobada de historias de usuario, desarrollo iterativo,
pruebas, documentación, capacitación y puesta en marcha.

\hypertarget{4-alcance-inicial}{%
\subsection{4. Alcance inicial}\label{4-alcance-inicial}}

La línea base funcional se consolidará y aprobará durante el
levantamiento. Se consideran inicialmente:

\begin{enumerate}
\def\labelenumi{\arabic{enumi}.}
\tightlist
\item
  identificación de usuarios, centros, roles y permisos;
\item
  inventario centralizado y ciclo de vida de los estudios;
\item
  registro y actualización de protocolos, responsables y criterios;
\item
  criterios de aceptación, selección y manejo de candidatos,
  seguimiento;
\item
  gestión de cupos para estudios, listas de espera, exámenes previos y
  alertas;
\item
  panel de indicadores de reclutamiento, compromisos y actividad;
\item
  trazabilidad de acciones relevantes;
\item
  datos de demostración, documentación y guía de uso.
\end{enumerate}

Las funciones de inteligencia artificial, procesamiento de voz,
integración con fichas clínicas y automatización de decisiones clínicas
se evaluarán durante el levantamiento. No se considerarán incorporadas a
la línea base hasta que exista una definición escrita de alcance,
seguridad, privacidad y criterios de aceptación.

La plataforma servirá como apoyo informativo y operativo. No determinará
por sí sola la elegibilidad clínica ni sustituirá el juicio profesional.

La interoperabilidad describe una capacidad de evolución del sistema. No
implica que el MVP incluya integración con ficha clínica u otros
sistemas que no hayan sido definidos, autorizados y acordados por
escrito.

\hypertarget{5-plan-de-trabajo}{%
\subsection{5. Plan de trabajo}\label{5-plan-de-trabajo}}

\hypertarget{mes-1--levantamiento-de-informaciuxf3n-requerimientos-y-flujo-de-procesosarquitectura}{%
\subsubsection{Mes 1 --- Levantamiento de información, requerimientos y
flujo de
procesos/arquitectura}\label{mes-1--levantamiento-de-informaciuxf3n-requerimientos-y-flujo-de-procesosarquitectura}}

\begin{itemize}
\tightlist
\item
  levantamiento y validación de procesos;
\item
  consolidación de requisitos e historias de usuario;
\item
  definición de arquitectura, modelo de datos y permisos;
\item
  priorización y aprobación de la línea base;
\item
  definición del flujo de procesos y la arquitectura inicial.
\end{itemize}

\hypertarget{mes-2--desarrollo-del-incremento-intermedio-del-mvp-y-validaciuxf3n-con-usuarios}{%
\subsubsection{Mes 2 --- Desarrollo del incremento intermedio del MVP y
validación con
usuarios}\label{mes-2--desarrollo-del-incremento-intermedio-del-mvp-y-validaciuxf3n-con-usuarios}}

\begin{itemize}
\tightlist
\item
  autenticación y autorización;
\item
  inventario y ciclo de vida de estudios;
\item
  criterios de aceptación, selección y manejo de candidatos;
\item
  pruebas continuas y validación con usuarios.
\end{itemize}

\hypertarget{mes-3--entrega-del-producto-operativo-puesta-en-marcha-y-marcha-blanca}{%
\subsubsection{Mes 3 --- Entrega del producto operativo, puesta en
marcha y marcha
blanca}\label{mes-3--entrega-del-producto-operativo-puesta-en-marcha-y-marcha-blanca}}

\begin{itemize}
\tightlist
\item
  historias de usuario incluidas en la línea base aprobada;
\item
  indicadores, alertas y trazabilidad acordados;
\item
  marcha blanca y pruebas de aceptación;
\item
  correcciones, documentación y capacitación;
\item
  puesta en marcha y entrega.
\end{itemize}

El MVP del mes 2 es un incremento intermedio; la entrega comprometida al
cierre del mes 3 es el producto operativo definido por la línea base
aprobada.

La gestión del proyecto seguirá principios compatibles con PMBOK 7,
adaptados a un desarrollo iterativo con entregas demostrables.

\hypertarget{6-precio-y-forma-de-pago}{%
\subsection{6. Precio y forma de pago}\label{6-precio-y-forma-de-pago}}

\begin{longtable}[]{@{}lrrr@{}}
\toprule\noalign{}
Concepto & Cantidad & Valor mensual neto & Total neto \\
\midrule\noalign{}
\endhead
\bottomrule\noalign{}
\endlastfoot
Desarrollo e implementación por hitos & 3 hitos & \$3.000.000 CLP netos
& \$9.000.000 CLP netos \\
IVA & & & \$1.710.000 CLP \\
\textbf{Total con IVA} & & & \textbf{\$10.710.000 CLP} \\
\end{longtable}

Adicionalmente, CICUC dispondrá de una bolsa de \textbf{USD 500 para
servicios AWS durante el desarrollo}. Este monto corresponde a costos de
infraestructura del proyecto, es de cargo de CICUC y no forma parte del
precio de ADYAC ni del total en pesos indicado arriba. La contratación,
cuenta y forma de pago de AWS deberán ser definidas por CICUC. Cualquier
ampliación de esta bolsa requerirá autorización escrita de CICUC antes
de incurrir en el gasto.

\begin{itemize}
\tightlist
\item
  Tres cuotas de \textbf{\$3.570.000 CLP IVA incluido}, facturadas
  contra la aceptación escrita de los hitos correspondientes a los meses
  1, 2 y 3.
\item
  El pago deberá efectuarse dentro de 15 días corridos desde la emisión
  del documento tributario respectivo.
\item
  Inicio estimado: sujeto a aceptación, disponibilidad de antecedentes y
  formalización administrativa.
\end{itemize}

Podrá acordarse por escrito una extensión de hasta un mes frente a
contingencias. La extensión no se presume gratuita: plazo, alcance y
contraprestación deberán ser negociados y aprobados por ambas partes
antes de iniciarla.

\hypertarget{7-entregables-y-aceptaciuxf3n}{%
\subsection{7. Entregables y
aceptación}\label{7-entregables-y-aceptaciuxf3n}}

\begin{itemize}
\tightlist
\item
  repositorios y código fuente del producto;
\item
  aplicación web operativa en el ambiente acordado;
\item
  base de datos y migraciones reproducibles;
\item
  pruebas correspondientes al alcance;
\item
  documentación técnica, operativa y de usuario;
\item
  capacitación inicial;
\item
  registro de historias de usuario y evidencia de aceptación.
\end{itemize}

Cada incremento se revisará contra criterios de aceptación escritos. El
cierre requiere una prueba de aceptación de la línea base aprobada. Las
observaciones deberán identificar el requisito o criterio incumplido.

CICUC dispondrá de 5 días hábiles para revisar cada entrega. La
aceptación debe constar por escrito y el silencio no constituye
aceptación. El procedimiento completo se encuentra en el Anexo A.

\hypertarget{8-propiedad-intelectual}{%
\subsection{8. Propiedad intelectual}\label{8-propiedad-intelectual}}

Adyac entregará a CICUC el código fuente, documentación y
configuraciones desarrolladas específicamente y pagadas en virtud del
proyecto, cediendo los derechos patrimoniales correspondientes en los
términos del Anexo C. CICUC podrá usar, reproducir, adaptar y modificar
estos entregables internamente o mediante otro proveedor para fines
institucionales.

Adyac conservará la titularidad de herramientas y componentes
preexistentes o desarrollados independientemente, que deberán
identificarse antes de su uso. Cuando sean indispensables para operar o
mantener el producto, otorgará a CICUC una licencia perpetua,
irrevocable, mundial y libre de regalías, suficiente para operar,
mantener y modificar el sistema, conforme al Anexo C. No se incluyen
componentes de machine learning en el alcance inicial.

Esta cláusula debe incorporarse al contrato definitivo y ser revisada
por las partes antes de su firma.

\hypertarget{9-control-de-cambios}{%
\subsection{9. Control de cambios}\label{9-control-de-cambios}}

La expresión ``historias de usuario faltantes'' se entenderá limitada a
las historias incluidas en la línea base aprobada al finalizar el
levantamiento. Cualquier requisito nuevo o cambio material deberá
registrar impacto en plazo, costo, seguridad y pruebas, y requerirá
aprobación escrita antes de su implementación.

\hypertarget{10-supuestos-responsabilidades-y-exclusiones}{%
\subsection{10. Supuestos, responsabilidades y
exclusiones}\label{10-supuestos-responsabilidades-y-exclusiones}}

CICUC designará responsables disponibles para validar procesos,
contenidos y entregas, y proporcionará oportunamente documentación y
datos ficticios o debidamente autorizados.

No se tratarán datos personales o clínicos reales hasta que el Anexo B
haya sido aprobado y las medidas técnicas, organizativas y de seguridad
allí establecidas se encuentren implementadas y verificadas. Durante
desarrollo y pruebas se utilizarán preferentemente datos sintéticos o
correctamente anonimizados. La infraestructura productiva, integraciones
institucionales y servicios de terceros se cotizarán cuando se conozcan
sus requisitos.

No se incluyen costos de infraestructura cloud que excedan la bolsa AWS
de USD 500 indicada en la sección 6, licencias de terceros,
integraciones no identificadas en el Anexo A, soporte adicional al
incluido en el Anexo A ni operación 24/7.

\hypertarget{11-aceptaciuxf3n}{%
\subsection{11. Aceptación}\label{11-aceptaciuxf3n}}

La aceptación deberá constar por escrito y comprender conjuntamente esta
cotización y los Anexos A, B y C, versión 1.0 de 31 de julio de 2026. La
persona que acepte por cada parte deberá contar con facultades
suficientes. Una orden de compra debe identificar expresamente los
cuatro documentos.

\begin{longtable}[]{@{}ll@{}}
\toprule\noalign{}
Por Adyac & Por CICUC \\
\midrule\noalign{}
\endhead
\bottomrule\noalign{}
\endlastfoot
Investigación, Tecnología y Gestión ADYAC SpA & \texttt{{[}NOMBRE{]}} \\
\texttt{{[}CARGO{]}} & \texttt{{[}CARGO{]}} \\
Firma: \_\_\_\_\_\_\_\_\_\_\_\_\_\_\_\_\_\_\_\_ & Firma:
\_\_\_\_\_\_\_\_\_\_\_\_\_\_\_\_\_\_\_\_ \\
Fecha: \_\_\_\_\_\_\_\_\_\_\_\_\_\_\_\_\_\_\_\_ & Fecha:
\_\_\_\_\_\_\_\_\_\_\_\_\_\_\_\_\_\_\_\_ \\
\end{longtable}

\end{document}
